Pembangunan Jalan Tol Yogyakarta--Solo berpotensi mengubah pola pergerakan pada jaringan jalan permukaan di wilayah Yogyakarta, termasuk koridor Ring Road Utara (RRU). Untuk memahami kondisi dasar sebelum perubahan jaringan tersebut dievaluasi lebih lanjut, diperlukan karakterisasi \textit{baseline} spatiotemporal berbasis data yang mampu mencakup periode dan area pengamatan luas. Penelitian ini menganalisis perjalanan kendaraan terindikasi pada enam persimpangan utama RRU menggunakan \textit{Mobile Positioning Data} (MPD) aktif. Dataset awal wilayah kajian terdiri dari 15.386.869 titik GPS dari 370.778 MAID. Setelah penyaringan spasial berbasis R-tree, \textit{bilocation collapse}, klasifikasi indikatif moda kendaraan bermotor, \textit{map matching}, dan segmentasi trajektori, diperoleh 212.407 ping dari 8.382 MAID dalam 38.973 trip valid. Analisis difokuskan pada matriks OD zona persimpangan, indikator intensitas persimpangan, distribusi ping per trip, dan pola OD zona eksploratif. Hasil mengindikasikan bahwa 28.767 trip memiliki pasangan OD zona teridentifikasi, dengan 58,48\% berada pada diagonal utama atau O=D. Pasangan OD non-diagonal terbesar adalah Condongcatur--UPN sebanyak 956 trip. Indikator intensitas menunjukkan Condongcatur dan Monjali sebagai zona dengan rata-rata MAID unik harian tertinggi, masing-masing 39,6 dan 39,3 MAID/hari. Distribusi trajektori menunjukkan bahwa data bersifat sparse, dengan median 4 ping per trip dan 70,37\% trip memiliki paling banyak 5 ping. Oleh karena itu, hasil penelitian ditafsirkan sebagai indikator kendaraan teramati berbasis sampel, bukan volume lalu lintas aktual.

\textbf{Kata kunci:} \textit{mobile positioning data}, Ring Road Utara, \textit{origin-destination}, intensitas persimpangan, trajektori sparse
